\section{Results}
\label{sec:results}

\subsection{Placement on the block}
\label{subsec:block}

\begin{figure*}[t]
  \centering
  \includegraphics[width=0.88\textwidth]{f2b-crossover-vs-d.png}
  \caption{Generation time under splitting divided by generation time under replication on
    the block benchmark, against the number of devices. Below 1, splitting is faster; above
    1, replication is faster. Panels separate platform and matrix width; marker shape gives
    the population. Each bar spans the smallest and largest ratio that any pair of repeats
    gives (the fastest split repeat over the slowest replicated one, up to the slowest split
    over the fastest replicated), across \BlockRepeats{} repeats per placement; the bars are
    not confidence intervals. Shown are dense, seed and rank 1;
    Appendix~\ref{sec:grids} has the complete grids with every variant of
    Table~\ref{tab:arms}.}
  \label{fig:f2}
\end{figure*}

At eight devices in one host, splitting is faster for every full-rank configuration. For
dense and seed, the variants in Figure~\ref{fig:f2}, splitting makes a generation
\FullRankSpeedup$\times$ faster. Dense alone gains \DenseSpeedupGPU$\times$ on the A100 and
\DenseSpeedupTPU$\times$ on the v5e, so the advantage does not come from seed's slower,
one-at-a-time evaluation. For rank-1 perturbations at width 2048, replication is faster:
splitting takes \LowRankSlowdownGPU\% longer on the A100 and \LowRankSlowdownTPU\% longer on
the v5e. These low-rank ranges cover both mirrored sampling, which Figure~\ref{fig:f2}
shows, and non-mirrored sampling; Appendix~\ref{sec:grids} gives every variant and the
absolute times.

We call a comparison unresolved when its repeat range, calculated as in
Figure~\ref{fig:f2}, contains 1. At eight devices, \NarrowLowRankUnresolved{} of the \NarrowLowRank{} low-rank configurations at width
512 are unresolved: the two placements take about the same time. Every full-rank
configuration, and every low-rank configuration at width 2048, is resolved. The repeats come
from one compiled program in one session, so they show variation within a run, not between
sessions or machines.

% generated by experiments/phase2/paper_evidence.py; do not edit
\begin{table*}[t]
\centering
\small
\caption{Worked example of Eq.~\ref{eq:crossover} at width 2048, population 256 and eight devices; all times in milliseconds. $C_A$ and $C_B$ are the separately measured reconstruction times. $T_{\mathrm{add}}$ is the time the all-reduce adds to the split reconstruction, measured in context; it stands in for $T_{\mathrm{ar}}(4P)$. The expected difference is $C_A-C_B+T_{\mathrm{ag}}-T_{\mathrm{add}}$, where $T_{\mathrm{ag}}$ is replication's second coefficient gather. Positive differences favor splitting. $t_A$ is the replicated generation time, the denominator of the ratio in Figure~\ref{fig:f2}.}
\label{tab:worked}
\begin{tabular}{llrrrrrr}
\toprule
platform & variant & $C_A$ & $C_B$ & $T_{\mathrm{add}}$ & expected $t_A-t_B$ & measured $t_A-t_B$ & $t_A$ \\
\midrule
A100 & seed & 61.05 & 7.78 & 1.25 & 52.03 & 54.30 & 113.26 \\
A100 & rank 1 & 0.77 & 0.35 & 1.62 & -1.19 & -1.40 & 19.88 \\
v5e & seed & 233.29 & 29.51 & 2.03 & 201.75 & 213.30 & 467.58 \\
v5e & rank 1 & 0.55 & 0.33 & 2.36 & -2.14 & -2.04 & 9.73 \\
\bottomrule
\end{tabular}
\end{table*}


Table~\ref{tab:worked} evaluates Eq.~\ref{eq:crossover} at one configuration on each
platform. For seed, splitting saves \ExSeedSaving\,ms of reconstruction and adds
\ExSeedComm\,ms of communication, so it is faster. For rank 1, reconstruction takes under
1\,ms under either placement, and the added communication outweighs the saving:
By Eq.~\ref{eq:crossover}, splitting is expected to be \ExRankPredGPU\,ms slower on the A100
and \ExRankPredTPU\,ms slower on the v5e, against \ExRankMeasGPU\,ms and \ExRankMeasTPU\,ms
measured. Relative to the replicated generation, the measured penalty is
\ExRankPenaltyGPU\% on the A100 and \ExRankPenaltyTPU\% on the v5e, where the generation is
shorter (\ExRankGenTPU\,ms against \ExRankGenGPU\,ms).

Figure~\ref{fig:f2}'s ratio and Eq.~\ref{eq:crossover}'s difference describe the same
quantity: $t_B/t_A=1-(t_A-t_B)/t_A$. A penalty of the same size in milliseconds is therefore
a larger ratio when the generation is shorter. Table~\ref{tab:tb8} in
Appendix~\ref{sec:timemodel} gives the all-reduce times measured in isolation; those exclude
any scheduling changes the all-reduce causes inside the compiled program, which
$T_{\mathrm{add}}$ in Table~\ref{tab:worked} includes.

\begin{figure*}[t]
  \centering
  \includegraphics[width=\textwidth]{f9-e17-crossover.png}
  \caption{Generation time under splitting divided by generation time under replication for
    Qwen2.5-0.5B on the TPU v5e, next-token loss evaluation included. Below 1, splitting is
    faster. The strips below the plots mark configurations that ran, ran out of memory
    (OOM), or were excluded; gaps are not connected. Only records with a fully known code version
    are used, and one-device records are excluded. Table~\ref{tab:tb5} gives the absolute
    times.}
  \label{fig:f9}
\end{figure*}

With separately measured reconstruction and communication times, Eq.~\ref{eq:crossover}
gives the faster placement in all \ModelResolved{} resolved configurations and in
\ModelCorrect{} of all \ModelConfigs{}; its \ModelMisses{} misses are all unresolved. This
tests whether separately timed components account for complete generations; nothing is
fitted to the measured differences. A simpler version that assumes $C_B=C_A/D$ and uses the
isolated all-reduce times misses \IdealMissesGPU{} resolved A100 configurations.
Appendix~\ref{sec:timemodel} gives both comparisons and the size of the remaining errors.

\subsection{Placement on a real model}



On Qwen2.5-0.5B on the TPU, the faster placement is the same as on the block
(Figure~\ref{fig:f9}). Only records whose exact code version is known are used, which
leaves \QwenFullComparisons{} full-rank and \QwenLowComparisons{} low-rank
comparisons; every eight-device record qualifies. Splitting is faster for full rank, by more
at eight devices than at four. Replication is faster in every low-rank comparison. The
closest, rank 16 at $N=128$, has a ratio of \QwenNearTieRatio{} and a repeat range starting at
\QwenNearTieBar{}, so it is resolved, but only just. Appendix~\ref{sec:memory} lists the
excluded records; they were excluded because their code version is not fully known, not
because they ran out of memory.

The update has \QwenParams{} float32 values (the embedding, which the output layer shares,
counted once), or \QwenUpdateGiB\,GiB. The all-reduce rate measured on the v5e at 96\,MiB,
scaled linearly to that size, gives \QwenCommEstimate\,ms; this extrapolates the measured
rate and fits nothing to Qwen's timings. In the four- and eight-device comparisons,
splitting adds \QwenLowPenalty\,ms for ranks 1 and 4. The all-reduce of the larger update
therefore accounts for the size of the penalty, although this neither separates the
reconstruction's share nor measures the all-reduce at \QwenUpdateGiB\,GiB directly.

At eight devices, rank 1's penalty stays between \QwenRankOnePenaltyMin{} and
\QwenRankOnePenaltyMax\,ms as the population grows from 32 to 128, while evaluation takes
longer, so its ratio falls from \QwenRankOneRatioSmallN{} to \QwenRankOneRatioLargeN{}. Adding
devices shortens evaluation and, in this comparison, also increases the absolute penalty:
at $N=32$, rank 1's ratio rises from \QwenRankOneRatioFourDevices{} on four devices to
\QwenRankOneRatioSmallN{} on eight. Rank 16's penalty falls as the population grows, to
\QwenRankSixteenPenalty\,ms at $N=128$. Before running this grid, predictions were made based
on an earlier study of two variants. Two held: which placement is faster for each variant,
and a smaller margin for rank 16 than for rank 1. The third did not: the margin was
predicted to shrink as rank grows, but at $N=64$ rank 4's ratio is
\QwenRankFourRatioMidN{}, above rank 1's \QwenRankOneRatioMidN{}. The population trend was
not predicted.

At eight devices and $N=240$, every low-rank variant runs out of memory under both
placements, while full rank fits. Low-rank candidates are evaluated together in one batch,
which is faster but holds every candidate's activations at once; seed evaluates candidates
one at a time and holds only one candidate's activations and noise. Compilation probes on a
reduced model support this explanation (Appendix~\ref{sec:memory}), but they do not show
where Qwen's memory limit would fall if both perturbation types were evaluated the same
way.

\subsection{Across a slow connection between hosts}
\label{subsec:boundary}

% generated by experiments/phase2/multihost/tb7_e18.py; do not edit
\begin{table}[t]
\centering
\small
\setlength{\tabcolsep}{3pt}
\caption{Generation time under replication minus generation time under splitting, $t_A-t_B$, in ms; positive values favor splitting, as in Table~\ref{tab:worked}. Columns give hosts $\times$ devices per host; the two hosts are connected by TCP sockets. Table~\ref{tab:tb7-audit} sets them beside the values Eq.~\ref{eq:crossover} gives.}
\label{tab:tb7}
\begin{tabular}{llrrrr}
\toprule
variant & $d$ & $N$ & $1\times8$ & $2\times4$ & $2\times8$ \\
\midrule
seed & 2048 & 128 & $+21.4$ & $-86.8$ & $-74.1$ \\
seed & 2048 & 256 & $+50.8$ & $-57.2$ & $-56.7$ \\
seed & 512 & 1024 & $+20.1$ & $+13.0$ & $+24.8$ \\
rank 1 & 2048 & 128 & $-1.6$ & $-108.9$ & $-159.1$ \\
rank 1 & 2048 & 256 & $-1.4$ & $-111.4$ & $-130.9$ \\
rank 1 & 512 & 1024 & $0.0$ & $-6.9$ & $-7.8$ \\
\bottomrule
\end{tabular}
\end{table}


Communication can reverse the full-rank preference. \HostConfigsCap{} of the block
configurations were measured again on two A100 nodes connected by TCP sockets, where the
effective cross-host all-reduce rate was \HostRate\,GiB/s.

With the same eight devices divided between two hosts (the $1\times8$ and $2\times4$ columns
of Table~\ref{tab:tb7}), $t_A-t_B$ falls by \HostBoundaryDrop\,ms at width 2048 for both
variants. For seed, that exceeds the \HostSeedSaving\,ms that splitting saved on one host, so
replication becomes faster. At width 512 the update is 6\,MiB, and splitting remains faster
for seed. For rank 1 at width 2048, replication's existing advantage grows across the
slow connection; at width 512, the near-tie becomes a clear advantage for replication.

Eq.~\ref{eq:crossover} also gives the connection speed below which splitting stops being
faster. For seed at width 2048 on eight devices, and ignoring latency, the 96\,MiB all-reduce
must run at \HostBreakEvenSmallN\,GiB/s or more at $N=128$, and at \HostBreakEvenLargeN\,GiB/s
or more at $N=256$, for splitting to remain faster. At those rates the all-reduce takes
exactly as long as the one-host saving plus the one-host all-reduce it replaces. These are
effective all-reduce rates, not network line rates; Appendix~\ref{sec:boundary-audit} gives
the calculation and the bounds from the repeats.

With the all-reduce rates measured on these hosts, the values Eq.~\ref{eq:crossover} gives
(expected, in Table~\ref{tab:tb7-audit}) match the faster placement in all
\HostExpectedMatches{} configurations, in both two-host layouts. On eight devices they differ
from the measurements by at most \HostExpectedErrorEight\,ms; on sixteen they overestimate the
penalty for splitting by up to \HostExpectedOverSixteen\,ms. This experiment shows what a slow
connection does to the choice of placement; it does not measure a high-bandwidth multi-host
cluster.

\subsection{What low-rank perturbations save}

\begin{figure*}[t]
  \centering
  \includegraphics[width=0.9\textwidth]{f4-cost-common-bfloat16.png}
  \caption{Generation time relative to dense (stored full-rank noise) for seed (left) and
    rank 1 (right), on one device with bfloat16 storage and default matrix-multiplication
    precision. Each row is one model width and population, measured on both platforms, at
    which every compared variant fits. Values below 1 are faster than dense.
    Figure~\ref{fig:f4-full} has the full grid, including memory limits.}
  \label{fig:f4}
\end{figure*}

On the block, rank 1 saves more time relative to dense on the v5e than on the A100
(Figure~\ref{fig:f4}). Across the \CostCommonConfigs{} configurations at which every compared
variant fits on both platforms, rank-1 generation time is \RankOneOverDenseGPU{} times the
dense time on the A100 and \RankOneOverDenseTPU{} times on the v5e (geometric means).
Regenerating full-rank noise trades speed for memory: seed is slower than dense at every
configuration where both were measured, but it fits every configuration of the cost grid,
unlike dense and rank 1. Appendix~\ref{sec:systems-ablations} has the full grid and the
precision comparisons. The difference between the platforms persists with matched precision: in
the single-device scaling sweeps, which use float32 storage and the highest precision on both
platforms, rank 1 takes \MatchedRankOneGPU{} times the dense time on the A100 and
\MatchedRankOneTPU{} times on the v5e, as geometric means over \MatchedPairs{} matched
configuration--placement pairs. These ratios describe the tested workload and implementation;
they do not identify the hardware mechanism behind the difference.

As an implementation check, we compared the library's single-device rank-1 throughput
with EGGROLL's released code: it is
\RefRatioGPU{} times EGGROLL's on the A100 and \RefRatioTPU{} times on the v5e. The cause
of the v5e gap was not found. Appendix~\ref{sec:refimpl} also gives the library's throughput
on eight devices, for which the single-device reference code has no counterpart.
